% Newton or Lagrange for the quasi-steady trim -- why the Lagrangian form is
% not faster. Markdown twin: lagrangian_formulation.md. Numbers from
% compare_lagrangian_residual.py (available on request)
% Build: latexmk -pdf lagrangian_formulation.tex
\documentclass[11pt,a4paper]{article}
\usepackage[margin=2.4cm]{geometry}
\usepackage{amsmath,amssymb}
\usepackage{booktabs}
\usepackage{siunitx}
\usepackage{xurl}
\usepackage[hidelinks]{hyperref}
\newcommand{\code}[1]{\texttt{#1}}
\newcommand{\vect}[1]{\boldsymbol{#1}}
\sisetup{per-mode=symbol}

\title{Newton or Lagrange for the quasi-steady trim:\\ why the Lagrangian form is not faster}
\author{AWETrim notes}
\date{10 October 2026}

\begin{document}
\maketitle

\paragraph{Question.} The reduced-order model writes the kite's dynamics as a
force balance in the course-aligned reference frame
(\code{SystemModel.force\_residual}, $-m\vect a + \vect F_\mathrm{ext}$; Cayon,
van Deursen, Schmehl 2026, WES 11, 1097). Would deriving the same model from a
Lagrangian make the trajectory NLP cheaper?

\paragraph{Answer.} No. The Lagrangian derivation yields the same three
equations, the inertia terms are \SI{0.3}{\percent} of each node function, and
function evaluations are a quarter of IPOPT's wall time. The numbers below come
from \path{compare_lagrangian_residual.py} (available on request)
(LEI-V3, periodic B-spline reel-out, $M=10$, flight-corrected aerostructural
ROM).

\section{The two derivations}

The path is prescribed on the unit sphere, $\vect q(s) = (\phi(s), \beta(s))$,
and the kite sits at $\vect R = r\,\vect q(s)$. The course frame has $x$ along
the path tangent, $z$ radially outward and $y$ normal to the path on the sphere.

\paragraph{Newton (library).} Velocity $\vect v = (v_\tau, 0, v_r)$ with
$v_\tau = r\,\sigma'\,\dot s$, $\sigma' = |\vect q'|$, and the frame rotation
rate
\begin{equation}
  \vect\omega = \Bigl(0,\ \frac{v_\tau}{r},\ \frac{v_\tau}{r}\tan\beta\sin\chi - \dot\chi\Bigr).
\end{equation}
The acceleration is $\vect a = (\dot v_\tau, 0, \dot v_r) + \vect\omega\times\vect v$;
the quasi-steady model drops the first bracket. Three rows:
$\vect F_\mathrm{ext} - m\vect a = \vect 0$.

\paragraph{Lagrange.} Generalised coordinates $(s, r)$, the path being a
holonomic constraint. Kinetic energy
$T = \tfrac{m}{2}\bigl(\dot r^2 + r^2\sigma'^2\dot s^2\bigr)$, generalised
forces $Q_s = \vect F\cdot\partial\vect R/\partial s = r\sigma' F_x$ and
$Q_r = F_z$. The Euler--Lagrange equations are
\begin{align}
  s:&\quad \frac{d}{dt}\bigl(m r^2\sigma'^2\dot s\bigr) - m r^2\sigma'\sigma''\dot s^2 = r\sigma' F_x, \\
  r:&\quad m\bigl(\ddot r - r\sigma'^2\dot s^2\bigr) = F_z .
\end{align}
Expanding the $s$ row gives $m r\sigma'\,(\dot v_\tau + v_\tau v_r / r)$, so the
quasi-steady assumption $\dot v_\tau = 0$ is a condition on $\ddot s$, not
$\ddot s = 0$. The lateral force the steering has to supply is the constraint
multiplier of the path, $m k_g v_\tau^2 / r$, with $k_g$ the geodesic
curvature. Divided by the metric factors, the three rows are
\begin{equation}
  F_x - m\frac{v_\tau v_r}{r} = 0, \qquad
  F_y + m k_g\frac{v_\tau^2}{r} = 0, \qquad
  F_z + m\frac{v_\tau^2}{r} = 0,
\end{equation}
which are the Newton rows. Over 200 random states the two residuals agree to
$3.5\times10^{-6}$ relative; the difference is the \code{fmax} floor on
$v_\tau$ and the $10^{-6}$ guard in $\tan\beta$ of the library form. The
multiplier row cannot be dropped, since it is the equation that determines
$u_s$, so no equation and no unknown is eliminated.

\section{Where the operations are}

\begin{table}[h]
\centering
\caption{SX instruction counts of the per-node trim residual.}
\begin{tabular}{lrr}
\toprule
expression & full spline & local support (NLP) \\
\midrule
Newton residual & 9971 & 3429 \\
Lagrange residual & 20585 & 6463 \\
path kinematics $\phi, \beta, \chi, \dot\chi, v_\tau$ & 9195 & \\
aerodynamic force (ROM) & 9846 & \\
inertia terms given the kinematics (Newton) & 26 & \\
\bottomrule
\end{tabular}
\end{table}

The whole acceleration costs 8877 instructions, but 26 of them are the inertia;
the rest is the spline evaluation and its $s$-derivatives, which the
aerodynamic force needs as well (the apparent wind is built in the course
frame). The Lagrangian residual, written through \code{curvature\_geodesic}, is
twice as large because the derivative chain of the spline is built a second
time. The Jacobians scale the same way (13460 against 22074).

\section{Where the time is}

Reel-out NLP, 120 nodes, $C_\phi$, $C_\beta$ and the depower free, the
production IPOPT options (\code{Phase.run\_opti}: L-BFGS, MUMPS, adaptive
barrier):

\begin{center}
\begin{tabular}{lr}
\toprule
iterations & 56 \\
IPOPT wall time & \SI{1.4}{s} \\
function evaluations & \SI{0.36}{s} (\SI{26}{\percent}), \SI{0.30}{s} of it the inequality Jacobian \\
MUMPS back-solves and barrier updates & \SI{74}{\percent} \\
inertia terms' share of the function evaluations & about \SI{0.3}{\percent}, $\sim$\SI{1}{ms} \\
\bottomrule
\end{tabular}
\end{center}

A reformulation of the inertia terms can therefore change at most a
millisecond per solve. The levers that remain are the spline-kinematics graph
in \code{ParametrizedKinematics} (nested \code{ca.gradient} chains, about
\SI{90}{\percent} of the node operations), the iteration count, and the KKT
size.

\section{Follow-up: the remaining levers, measured}

Same 120-node reel-out NLP, single runs, so differences of a few milliseconds
per iteration are noise.

\paragraph{Spline-kinematics graph.} The properties of
\code{ParametrizedKinematics} rebuild the spline sub-graph on every access
($\phi$ alone is evaluated about nine times inside one node function) and the
SX expansion keeps the copies. Written with analytic B-spline derivatives, the
course-frame chain $\chi, v_\tau, \dot\chi, k_g$ costs 470 instructions
instead of 4951 (agreement $10^{-11}$). The cheapest fix, caching the
properties so each sub-expression is built once, gives most of that:

\begin{center}
\begin{tabular}{lrr}
\toprule
node function & properties & cached \\
\midrule
residual & 3429 & 2280 \\
angle of attack & 2868 & 1719 \\
curvature triple & 5626 & 2013 \\
\bottomrule
\end{tabular}
\end{center}

Yet the solve does not get faster: 56 iterations either way, 24 against
\SI{26}{ms} per iteration without turn-radius rows, 31 against \SI{28}{ms}
with them ($R_\mathrm{min} = \SI{25}{m}$, 95 against 94 iterations). Function
evaluations stay at \SI{25}{\percent} of the wall time; the rest is the KKT
back-solve, which L-BFGS turns into many right-hand sides per iteration. The
smaller graph still cuts the whole-graph expansion memory in proportion, which
matters at the full-cycle size (\SI{3.7}{MB} per node at $N = 376$), and would
be a safe change: the mathematics is unchanged.

\paragraph{Hessian.} The exact Hessian (available from the SX graph) halves the
iterations, 26 instead of 56, but the Hessian evaluation costs \SI{0.69}{s}
and the solve ends at \SI{1.66}{s} against \SI{1.36}{s}. With the turn-radius
rows it is far worse: 245 iterations and \SI{29}{s}, the $\sigma'^3$
polynomial rows giving an ill-conditioned Hessian. A 20-deep L-BFGS memory
changes nothing (57 and 84 iterations). The production choice, L-BFGS with the
default memory, stands.

\paragraph{Conclusion.} At this problem size the NLP is already balanced
between function evaluation and linear algebra, and no single lever gains
more than the \SI{25}{\percent} that function evaluation represents. Wall time
is governed by the iteration count and the number of staged passes, not by
the cost of the model.

\paragraph{Reproduce.} The A/B uses the flags \code{--turn-radius 25},
\code{--kinematics-cache} and \code{--exact-hessian} of
{\small
\begin{verbatim}
python compare_lagrangian_residual.py --nlp   (available on request)
\end{verbatim}
}

\end{document}
